\section{Definición del problema de las \texorpdfstring{$p$}{p}-medianas}
\label{sec:definicion}

\subsection{Datos, decisión y objetivo}
\begin{description}[leftmargin=1.5em]
  \item[Conjuntos.] Clientes $[N]=\{1,\dots,N\}$ y sitios candidatos $[M]=\{1,\dots,M\}$. En
        muchas instancias geométricas (k-medoids, TSP, BIRCH) clientes y sitios son el mismo
        conjunto de puntos, de modo que $N=M$.
  \item[Parámetros.] Distancias $d_{ij}\ge 0$ entre cliente $i$ y sitio $j$; entero $p$ con
        $1\le p\le M$ (número de sitios a abrir).
  \item[Decisión.] Elegir $S\subseteq[M]$ con $|S|=p$.
  \item[Objetivo.]
        \[
          \min_{\,|S|=p}\ \sum_{i=1}^{N}\ \min_{j\in S}\ d_{ij}.
        \]
\end{description}

Cada cliente se asigna a su \textbf{sitio abierto más cercano}; su \textbf{distancia de
asignación} es esa distancia mínima. La función objetivo es la suma de las distancias de
asignación de todos los clientes.

\subsection{Estructura que será explotada}
La observación central, que motiva tanto la relajación lagrangiana como Benders, es:

\begin{quote}
Si \emph{fijamos} qué sitios están abiertos, el problema se \textbf{separa por cliente} y se
vuelve trivial (cada cliente toma su mínimo). La dificultad reside únicamente en la elección
combinatoria del conjunto $S$.
\end{quote}

Las variables que indican qué sitios abrir son, por tanto, las \emph{variables complicantes};
toda la maquinaria de descomposición se construye alrededor de fijarlas o relajarlas. En el
lenguaje del curso (Cap.~3), las aperturas $y$ son las decisiones \textbf{estratégicas o de
diseño} (discretas), y la asignación de clientes es la decisión \textbf{operacional} (fácil dada
$y$); el acoplamiento entre ambas es lo que genera la dificultad, no la integralidad por sí sola.

\subsection{Caso base de verificación: \texttt{toy1}}
Para validar toda la cadena (parser, distancias, separación, fases) se usa una instancia mínima
\texttt{toy1} con $N=M=4$, $p=2$ y coordenadas 2D. Con distancia euclidiana truncada las
distancias son $d(0,1)=3,\ d(0,2)=4,\ d(0,3)=5,\ d(1,2)=5,\ d(1,3)=4,\ d(2,3)=3$. Por
enumeración exhaustiva el óptimo es $\mathbf{6}$ (abrir $\{0,2\}$). Este valor se reproduce por
fuerza bruta, por el modelo de referencia F3, por la Fase 1 y por la Fase 2 de la implementación
(archivo \texttt{instances/toy/toy1.pmp}; verificado por \texttt{make test}).
